%% chapters/09-eigenvalue-convergence.tex
\chapter{Eigenvalues Are Continuous Functions of the Matrix}
\label{ch:eigconv}

\whereWeAre{%
This is Joseph's Proposition~3.1: when matrices converge, their eigenvalues converge. In the deterministic, almost-sure, and in-probability flavours, the same continuity argument carries the day. This is the spine on which every later eigenvalue convergence claim hangs.%
}

\PLACEHOLDER{Sources: existing main.tex `Continuation: Sample Covariance Matrices, Norms, and Eigenvalue Convergence' section.}

\section{The logical chain}
\PLACEHOLDER{Determinant continuous (polynomial in entries) $\to$ characteristic polynomial coefficients continuous $\to$ roots continuous $\to$ eigenvalues continuous.}

\section{Three convergence flavours}
\PLACEHOLDER{Deterministic, a.s., in probability. Each follows from the deterministic version applied either pointwise on a measure-1 set or via continuity bounds.}

\section{Sanity check: the $2 \times 2$ case}
\PLACEHOLDER{Explicit eigenvalue formula in 2D; entrywise convergence trivially gives eigenvalue convergence.}

\section{Brief warm-ups}
\PLACEHOLDER{3--5 short questions.}

\chapterRecap{%
\begin{itemize}
  \item \PLACEHOLDER{Matrix convergence implies eigenvalue convergence.}
  \item \PLACEHOLDER{The pipeline is: entries $\to$ coefficients $\to$ roots.}
  \item \PLACEHOLDER{This is the analytic engine for the next several chapters.}
\end{itemize}%
}

\section*{Exercises}
\PLACEHOLDER{8--15 longer exercises.}
